\section{Proofs of NTK Recursions}
\label{appendix:ntk-recursions}

\subsection{Proof of Theorem~\ref{thm:ntk_recursion}: Infinite-width NTK recursion}
\label{appendix:proof-ntk-recursion}

\begin{proof}
We prove the recursion by decomposing the NTK at layer \( \ell \) into contributions of the trainable parameters up to layer \( \ell \), and invoking the Law of Large Numbers (LLN) under the sequential infinite-width limit. Let \( f^{(\ell)} \) denote the network output truncated at layer \( \ell \). For finite width, the (random) NTK reads
\[
\Theta^{(\ell)}_{\mathrm{rand}}(x,x') \;=\; \underbrace{\big(\nabla_{\theta^{(\ell-1)}} f^{(\ell)}(x)\big)^\top \big(\nabla_{\theta^{(\ell-1)}} f^{(\ell)}(x')\big)}_{\text{Term 1}} \;+\; \underbrace{\sum_{j=1}^{n_\ell}\frac{\partial f^{(\ell)}(x)}{\partial A^{(\ell)}_j}\frac{\partial f^{(\ell)}(x')}{\partial A^{(\ell)}_j}}_{\text{Term 2}}.
\]

\paragraph{Term 2 (layer-\(\ell\) output weights \(A^{(\ell)}\)).}
By the forward pass, \(\partial f^{(\ell)}(x)/\partial A^{(\ell)}_j = \frac{\sigma_A}{\sqrt{n_\ell}}\ReLU\!\big(\bm{w}^{(\ell)\top}_j\bm{h}^{(\ell-1)}(x)+\beta b^{(\ell)}_j\big)\). Therefore
\[
\text{Term 2} \;=\; \frac{\sigma_A^2}{n_\ell}\sum_{j=1}^{n_\ell} \ReLU\!\big(\bm{w}^{(\ell)\top}_j\bm{h}^{(\ell-1)}(x)+\beta b^{(\ell)}_j\big)\,\ReLU\!\big(\bm{w}^{(\ell)\top}_j\bm{h}^{(\ell-1)}(x')+\beta b^{(\ell)}_j\big).
\]
Conditionally on \(\bm{h}^{(\ell-1)}\), the summands are i.i.d. over \((\bm{w}^{(\ell)}_j,b^{(\ell)}_j)\). As \(n_\ell\to\infty\) (with lower layers already taken to their limits), LLN yields
\[
\text{Term 2} \xrightarrow{p} \sigma_A^2\,\Sigma^{(\ell)}(x,x').
\]

\paragraph{Term 1 (previous layers).}
By the chain rule,
\[
\nabla_{\theta^{(\ell-1)}} f^{(\ell)}(x) \;=\; \Big(\frac{\partial f^{(\ell)}(x)}{\partial \bm{h}^{(\ell-1)}(x)}\Big)\,\nabla_{\theta^{(\ell-1)}} \bm{h}^{(\ell-1)}(x).
\]
Hence Term 1 factorizes as \(\Theta^{(\ell-1)}_{\mathrm{rand}}(x,x')\) multiplied by a Jacobian metric:
\[
\text{Term 1} \;=\; \Theta^{(\ell-1)}_{\mathrm{rand}}(x,x')\;\cdot\;\Big\langle \frac{\partial f^{(\ell)}}{\partial \bm{h}^{(\ell-1)}}(x), \frac{\partial f^{(\ell)}}{\partial \bm{h}^{(\ell-1)}}(x') \Big\rangle.
\]
Taking the inner product and applying LLN as \(n_\ell\to\infty\) yields
\[
\Big\langle \frac{\partial f^{(\ell)}}{\partial \bm{h}^{(\ell-1)}}(x), \frac{\partial f^{(\ell)}}{\partial \bm{h}^{(\ell-1)}}(x') \Big\rangle \xrightarrow{p} \sigma_A^2\,\dot{\Sigma}^{(\ell)}(x,x').
\]
Therefore,
\[
\text{Term 1} \xrightarrow{p} \Theta^{(\ell-1)}(x,x')\cdot \sigma_A^2\,\dot{\Sigma}^{(\ell)}(x,x').
\]

Combining the two limits completes the recursion
\[
\Theta^{(\ell)}(x,x') \;=\; \Theta^{(\ell-1)}(x,x')\cdot\sigma_A^2\,\dot{\Sigma}^{(\ell)}(x,x') \;+\; \sigma_A^2\,\Sigma^{(\ell)}(x,x').
\]
The sequential limit assumption ensures that lower-layer random objects have already converged when taking the next-layer limit, justifying the conditional LLN applications.
\end{proof}

\subsection{Proof of Two-layer NTK Corollary}
\label{appendix:proof-2layer}

\begin{proof}
Apply Theorem~\ref{thm:ntk_recursion} for \( \ell=1 \) and \( \ell=2 \). For \( \ell=1 \), since \( \Kop^{(0)}=0 \), the recursion gives:
\[
\Kop^{(1)}(x,x') = \Kop^{(0)}(x,x')\cdot\dot{\Sigma}^{(1)}(x,x') + \Sigma^{(1)}(x,x') = 0 + \Sigma^{(1)}(x,x') = \Sigma^{(1)}(x,x').
\]
For \( \ell=2 \), substitute \( \Kop^{(1)} = \Sigma^{(1)} \) into the recursion:
\[
\Kop^{(2)}(x,x') = \Sigma^{(1)}(x,x')\cdot\dot{\Sigma}^{(2)}(x,x') + \Sigma^{(2)}(x,x'),
\]
which is the stated expression.
\end{proof}

\subsection{Proof of General \texorpdfstring{\(L\)}{L}-layer NTK Corollary}
\label{appendix:proof-general-L}

\begin{proof}
We prove the explicit form by induction on \( L \).

\paragraph{Base case \( L=1 \).}
From Theorem~\ref{thm:ntk_recursion}, \( \Kop^{(1)} = \Sigma^{(1)} \). The formula with \( L=1 \) gives:
\[
\Kop^{(1)}(x,x') = \Sigma^{(1)}(x,x') \prod_{k=2}^{1} \dot{\Sigma}^{(k)}(x,x') = \Sigma^{(1)}(x,x') \cdot 1 = \Sigma^{(1)}(x,x'),
\]
where the empty product equals 1 by convention.

\paragraph{Inductive step.}
Assume the formula holds for depth \( L-1 \):
\[
\Kop^{(L-1)}(x,x') = \sum_{\ell=1}^{L-1} \Sigma^{(\ell)}(x,x') \prod_{k=\ell+1}^{L-1} \dot{\Sigma}^{(k)}(x,x').
\]
By Theorem~\ref{thm:ntk_recursion},
\[
\Kop^{(L)}(x,x') = \Kop^{(L-1)}(x,x')\cdot\dot{\Sigma}^{(L)}(x,x') + \Sigma^{(L)}(x,x').
\]
Substituting the inductive hypothesis:
\begin{align*}
\Kop^{(L)}(x,x') &= \left(\sum_{\ell=1}^{L-1} \Sigma^{(\ell)}(x,x') \prod_{k=\ell+1}^{L-1} \dot{\Sigma}^{(k)}(x,x')\right) \cdot \dot{\Sigma}^{(L)}(x,x') + \Sigma^{(L)}(x,x')\\
&= \sum_{\ell=1}^{L-1} \Sigma^{(\ell)}(x,x') \prod_{k=\ell+1}^{L} \dot{\Sigma}^{(k)}(x,x') + \Sigma^{(L)}(x,x'),
\end{align*}
where in the last step we extended the product to include \( \dot{\Sigma}^{(L)} \). The final term corresponds to \( \ell=L \) with an empty product. Therefore:
\[
\Kop^{(L)}(x,x') = \sum_{\ell=1}^{L} \Sigma^{(\ell)}(x,x') \prod_{k=\ell+1}^{L} \dot{\Sigma}^{(k)}(x,x'),
\]
completing the induction.
\end{proof}

\section{RKHS Theory and Concentration}
\label{appendix:rkhs-theory}

\subsection{No RKHS Advantage}
\label{appendix:no-rkhs-advantage}

\begin{theorem}[No advantage in RKHS]
\label{thm:appendix-no-rkhs-advantage}
Under the RF-LR setting with ReLU nonlinearity, isotropic random features on the sphere \( \mathbb{S}^{d-1} \), and in the kernel (NTK) regime, the depth-\(L\) NTK induces the same RKHS as the shallow (two-layer) ReLU kernel. In particular, the corresponding RKHSs coincide as sets with equivalent norms. We use this statement as provided by Theorem 1 in~\cite{bietti2021deepequalsshallowrelu}.
\end{theorem}

\begin{corollary}[Spectral decay on \( \mathbb{S}^{d-1} \) via endpoint expansions]
\label{cor:appendix-spectral-decay}
Let \( \kappa:[-1,1]\to\mathbb{R} \) be the rotation-invariant kernel associated to the (deep or shallow) ReLU NTK on \( \mathbb{S}^{d-1} \). Assume \( \kappa \in C^\infty(-1,1) \) and admits the asymptotic expansions, for \( t\ge 0 \),
\[
\kappa(1-t) = p_1(t) + c_1 t^{\nu} + o(t^{\nu}),\qquad
\kappa(-1+t) = p_{-1}(t) + c_{-1} t^{\nu} + o(t^{\nu}),
\]
where \( p_1,p_{-1} \) are polynomials, \( \nu>0 \) non-integer, and derivatives satisfy analogous expansions. Then there exists a constant \( C(d,\nu)>0 \) such that the eigenvalues \( \mu_k \) (in the spherical harmonics basis) satisfy:
\begin{itemize}
    \item For even \(k\) and \( c_1 \neq -c_{-1}\): \( \mu_k \sim (c_1+c_{-1})\,C(d,\nu)\,k^{-d-2\nu+1} \).
    \item For odd \(k\) and \( c_1 \neq c_{-1}\): \( \mu_k \sim (c_1-c_{-1})\,C(d,\nu)\,k^{-d-2\nu+1} \).
\end{itemize}
If \( |c_1|=|c_{-1}| \), then \( \mu_k = o\big(k^{-d-2\nu+1}\big) \) for one parity (or both if \(c_1=c_{-1}=0\)). If \( \kappa \) is \( C^\infty \) on \([{-}1,1]\) with no such \( \nu \), then \( \mu_k \) decays faster than any polynomial. See~\cite{bietti2021deepequalsshallowrelu}.
\end{corollary}

\subsection{Proof of Theorem~\ref{thm:rkhs_rank_concentration}: Rank-driven concentration}
\label{appendix:proof-rkhs-rank-concentration}

\begin{proof}
Let \( x_1,y_1\in\mathbb{R}^r \) be Gaussian projections with arbitrary fixed correlation \( \rho\in[-1,1] \). Define \( U=\|x_1\|/\sqrt{r} \), \( V=\|y_1\|/\sqrt{r} \), and \( W=UV=\|x_1\|\,\|y_1\|/r \). We prove that for \( \epsilon\in(0,1) \),
\[
\mathbb{P}\!\left(\left|W-1\right|\ge \epsilon\right)\le 4\exp\!\left(-\frac{r\epsilon^2}{8}\right),
\]
and that for \( \epsilon\ge 1 \) there exist constants \( c_1,c_2>0 \) with \( \mathbb{P}(|W-1|\ge \epsilon)\le c_1 e^{-c_2 r\epsilon} \).

The Euclidean norm is 1-Lipschitz, so by Gaussian concentration, for all \( \delta>0 \),
\[
\mathbb{P}\!\left(\left|U-1\right|\ge \delta\right)\le 2\exp\!\left(-\frac{r\delta^2}{2}\right),\qquad
\mathbb{P}\!\left(\left|V-1\right|\ge \delta\right)\le 2\exp\!\left(-\frac{r\delta^2}{2}\right).
\]
We use
\[
|UV-1|\;\le\;|U-1|+|V-1|+|U-1|\,|V-1|.
\]
Fix \( \epsilon\in(0,1) \) and set \( \delta=\epsilon/2 \). On the event \( \{|U-1|<\delta,\;|V-1|<\delta\} \), we have
\[
|UV-1|\;\le\;2\delta+\delta^2\;\le\; \epsilon.
\]
Therefore,
\[
\mathbb{P}\!\left(|UV-1|\ge \epsilon\right)\;\le\; \mathbb{P}\!\left(|U-1|\ge \delta\right)\;+\;\mathbb{P}\!\left(|V-1|\ge \delta\right) \;\le\;4\exp\!\left(-\frac{r\epsilon^2}{8}\right).
\]
For \( \epsilon\ge 1 \), standard tail bounds for \( \chi \)-norms (sub-exponential beyond the local regime) yield the \( \exp(-c r\epsilon) \) behavior.
\end{proof}

\section{Spectral Theory}
\label{appendix:spectral-theory}

\subsection{Asymptotic Signal-Noise Decomposition}
\label{appendix:signal-noise-decomposition}

\begin{proof}[Proof of Theorem~\ref{thm:kernel_clt}]
We establish the asymptotic distribution of the empirical kernel \( \hat{\Theta}^{(2)}_r = \Psi(\hat{\rho}_r,\hat{w}_r) \) via a joint CLT and the multivariate Delta method. Let \( \mathbf{Z}_r = (\hat{\rho}_r,\,\hat{w}_r)^\top \). For isotropic Gaussian projections in rank \( r \), Fisher's theory gives
\[
\sqrt{r}\,(\hat{\rho}_r-\rho) \;\xrightarrow{d}\; \mathcal{N}\!\left(0,\,(1-\rho^2)^2\right).
\]
Writing \( u=\|x\|^2 \), \( v=\|y\|^2 \) and \( \hat{w}_r=\sqrt{uv}/r \), Kibble's bivariate chi-square yields
\[
\sqrt{r}\,(\hat{w}_r-1) \;\xrightarrow{d}\; \mathcal{N}\!\left(0,\,1+\rho^2\right).
\]
Angular and radial components are independent for isotropic Gaussians, hence
\[
\sqrt{r}\,\Big(\mathbf{Z}_r - (\rho,1)^\top\Big) \;\xrightarrow{d}\; \mathcal{N}\!\Big(\mathbf{0},\ \mathbf{\Lambda}(\rho)\Big),
\quad
\mathbf{\Lambda}(\rho) \;=\; \begin{pmatrix}
(1-\rho^2)^2 & 0\\[2pt]
0 & 1+\rho^2
\end{pmatrix}.
\]
The map \( \Psi(u,w)=\Theta^{(1)}(u)\big(1-\arccos(u)/\pi\big) + 2w\,\Sigma^{(1)}(u)+1 \) is \( C^1 \) near \( (\rho,1) \). Therefore,
\[
\sqrt{r}\,\Big(\Psi(\mathbf{Z}_r) - \Psi(\rho,1)\Big) \;\xrightarrow{d}\; \mathcal{N}\!\Big(0,\ \nabla\Psi(\rho,1)^\top \mathbf{\Lambda}(\rho)\,\nabla\Psi(\rho,1)\Big).
\]
We have \( \partial_u \Psi(\rho,1) = \partial_\rho K_\infty(\rho) \) and \( \partial_w \Psi(\rho,1) = 2\,\Sigma^{(1)}(\rho) \). Using the diagonal covariance, we obtain
\[
\sigma_{\mathrm{tot}}^2(\rho)
\;=\; \big(\partial_\rho K_\infty(\rho)\big)^2\,(1-\rho^2)^2 \;+\; \big(2\,\Sigma^{(1)}(\rho)\big)^2\,(1+\rho^2).
\]
Hence
\[
\hat{\Theta}^{(2)}_r \;=\; K_\infty(\rho) \;+\; \frac{1}{\sqrt{r}}\,\mathcal{G}(\rho) \;+\; o_p(r^{-1/2}),
\quad
\mathcal{G}(\rho)\sim \mathcal{N}\!\big(0,\sigma_{\mathrm{tot}}^2(\rho)\big),
\]
which proves the claim.
\end{proof}

\subsection{Fisher and Kibble Distributions}
\label{appendix:fisher-kibble}

\begin{lemma}[Fisher–Kibble decoupling]
\label{lem:fisher-kibble}
Let \(x,y\) be input vectors and let \(x_1,y_1\) denote their rank-\(r\) random projections. Define the empirical correlation \( \rho_1 = \cos\angle(x_1,y_1) \) and squared norms \( u=\|x_1\|^2, v=\|y_1\|^2 \). Then:
\begin{itemize}
    \item \( \rho_1 \) follows Fisher's correlation distribution (1915), centered at the true correlation \( \rho \).
    \item \( (u,v) \) follow Kibble's (1941) bivariate Gamma (chi-square) law.
    \item Angular and radial parts are independent: \( p(\rho_1,u,v) = p_{\text{Fisher}}(\rho_1)\,p_{\text{Kibble}}(u,v) \).
\end{itemize}
\end{lemma}

\begin{proof}
Let \(x, y \in \mathbb{R}^d\) be fixed input vectors with correlation \(\rho = \langle x, y\rangle/(\|x\|\|y\|)\). Consider a rank-\(r\) random projection matrix \(P \in \mathbb{R}^{r \times d}\) with entries \(P_{ij} \sim \mathcal{N}(0, 1/d)\) i.i.d., so that \(x_1 = Px\) and \(y_1 = Py\) are the projected vectors.

Since \(P\) has i.i.d. Gaussian entries, the projected vectors \(x_1, y_1 \in \mathbb{R}^r\) form a \(\rho\)-correlated bivariate Gaussian pair. Specifically, for each component \(i = 1, \dots, r\), the pairs \((x_{1i}, y_{1i})\) are i.i.d. with joint distribution
\[
\begin{pmatrix} x_{1i} \\ y_{1i} \end{pmatrix} \sim \mathcal{N}\left(\mathbf{0}, \begin{pmatrix} \|x\|^2/d & \rho \|x\|\|y\|/d \\ \rho \|x\|\|y\|/d & \|y\|^2/d \end{pmatrix}\right).
\]
By rotational invariance of the Gaussian projection and the fact that the correlation \(\rho\) is preserved under orthogonal transformations, we can assume without loss of generality that the covariance structure factors into angular and radial components.

\paragraph{Angular part (Fisher distribution).}
The empirical correlation is defined as
\[
\rho_1 = \frac{\langle x_1, y_1 \rangle}{\|x_1\|\|y_1\|} = \frac{\sum_{i=1}^r x_{1i} y_{1i}}{\sqrt{\sum_{i=1}^r x_{1i}^2}\sqrt{\sum_{i=1}^r y_{1i}^2}}.
\]
This is the sample correlation coefficient of \(r\) pairs of correlated Gaussian random variables. Fisher (1915) showed that the distribution of \(\rho_1\) depends only on the true correlation \(\rho\) and the sample size \(r\), with density given in Theorem~\ref{thm:fisher-distribution}. The key observation is that \(\rho_1\) is a function purely of the angles between the projected vectors, independent of their magnitudes.

\paragraph{Radial part (Kibble distribution).}
The squared norms \(u = \|x_1\|^2 = \sum_{i=1}^r x_{1i}^2\) and \(v = \|y_1\|^2 = \sum_{i=1}^r y_{1i}^2\) are sums of correlated chi-square variables. Since \((x_{1i}, y_{1i})\) are \(\rho\)-correlated Gaussians, the bivariate distribution of \((u, v)\) follows Kibble's (1941) bivariate chi-square law. The joint density involves the modified Bessel function \(I_{\nu}(z)\), which appears because the dot product \(\langle x_1, y_1 \rangle = \sum_{i=1}^r x_{1i} y_{1i}\) can be expressed as a weighted sum of products of correlated normals, whose distribution relates to Bessel functions through the generating function of the bivariate chi-square distribution.

Specifically, the joint characteristic function of \((u, v)\) is
\[
\phi(s, t) = \mathbb{E}[e^{i(su + tv)}] = \left(1 - 2i(1-\rho^2)(s+t) - 4\rho^2 st\right)^{-r/2},
\]
and the inverse Fourier transform yields Kibble's density, which contains the modified Bessel function \(I_{\frac{r}{2}-1}\) as shown in Theorem~\ref{thm:kibble-distribution}. The Bessel function arises from the integral representation of the bivariate chi-square density when the correlation \(\rho \neq 0\).

\paragraph{Independence of angular and radial parts.}
Rotational invariance of isotropic Gaussian projections implies that \(\rho_1 = \cos\angle(x_1, y_1)\) and the norms \((\|x_1\|, \|y_1\|)\) are independent, yielding the factorization \(p(\rho_1, u, v) = p_{\text{Fisher}}(\rho_1) \cdot p_{\text{Kibble}}(u, v)\). This is the classical independence of the sample correlation from sample variances in bivariate normal data.
\end{proof}

\begin{theorem}[Fisher distribution of the sample correlation]
\label{thm:fisher-distribution}
Let \((X_i,Y_i)_{i=1}^r\) be i.i.d. centered Gaussian pairs with unit variances and correlation \(\rho\in[-1,1]\). Define the sample correlation
\[
r \;=\; \frac{\sum_{i=1}^r X_iY_i}{\sqrt{\sum_{i=1}^r X_i^2}\sqrt{\sum_{i=1}^r Y_i^2}} \;\in [-1,1].
\]
Then \(r\) has density (Fisher, 1915)
\[
p(r\mid \rho,r) \;=\; \frac{(r-2)\,\Gamma(r-1)\,(1-\rho^2)^{\frac{r-1}{2}}\,(1-r^2)^{\frac{r-4}{2}}}{\sqrt{2\pi}\,\Gamma\!\left(r-\tfrac{1}{2}\right)\,(1-\rho\,r)^{\,r-\frac{3}{2}}}\;\;{}_2F_1\!\left(\tfrac{1}{2},\tfrac{1}{2};\, r-\tfrac{1}{2};\, \tfrac{1+\rho\,r}{2}\right),
\]
for \(r>2\). For large \(r\), Fisher's \(z\)-transform \(z=\operatorname{arctanh}(r)\) satisfies
\[
z \;\sim\; \mathcal{N}\!\left(\operatorname{arctanh}(\rho)+\frac{\rho}{2(r-1)},\;\frac{1}{\,r-3\,}\right),
\]
so \(\operatorname{Var}(r)=O(1/r)\).
\end{theorem}

\begin{theorem}[Kibble bivariate chi-square law]
\label{thm:kibble-distribution}
Let \((X_i,Y_i)_{i=1}^r\) be as above and define the squared norms
\[
U=\sum_{i=1}^r X_i^2,\qquad V=\sum_{i=1}^r Y_i^2.
\]
Then \((U,V)\) has Kibble's (1941) bivariate chi-square density
\[
f(u,v) \;=\; \frac{(uv)^{\frac{r/2-1}{2}}\,\exp\!\left(-\tfrac{u+v}{2(1-\rho^2)}\right)}{\Gamma(r/2)\,\big(2(1-\rho^2)\big)^{\frac{r}{2}+1}\,\rho^{\frac{r/2-1}{2}}}\; I_{\,\frac{r}{2}-1}\!\left(\frac{\rho\sqrt{uv}}{\,1-\rho^2\,}\right),\qquad u,v\ge 0,
\]
where \(I_\nu\) is the modified Bessel function of the first kind. In particular,
\[
\mathbb{E}[U]=\mathbb{E}[V]=r,\quad \operatorname{Var}(U)=\operatorname{Var}(V)=2r,\quad \operatorname{Cov}(U,V)=2r\,\rho^2.
\]
\end{theorem}

\subsection{Proof of Theorem~\ref{thm:spike-bulk}: Deformed Marchenko–Pastur with Spike}
\label{appendix:marchenko-pastur}

\begin{proof}
We prove the spectral limit using Theorem 2.1 of~\cite{El_Karoui_2010} on the spectrum of inner product kernel random matrices. This connects the microscopic kernel structure (ReLU/Kibble calculations) with the macroscopic matrix behavior (spectrum).

\paragraph{Setting.}
Let \(\mathbf{M} = [\Theta^{(2)}_{ij}]_{i,j=1}^n\) be the \(n \times n\) NTK matrix with entries \(M_{ij} = \Theta^{(2)}(x_i, x_j) = K_\infty(\rho_{ij}) + O_p(r^{-1/2})\), where \(\rho_{ij} = \langle x_i, x_j \rangle / (\|x_i\|\|x_j\|)\) is the cosine similarity and \(K_\infty(\rho) = \Theta^{(1)}(\rho)(1 - \arccos(\rho)/\pi) + 2\Sigma^{(1)}(\rho) + 1\) is the deterministic two-layer kernel limit.

Under Assumption~\ref{assump:RMT}, we have \(n \sim r\) (so \(n/p\) remains bounded where \(p\) is the ambient dimension) and the data points \(x_i\) are normalized. The kernel function \(f(\rho) = K_\infty(\rho)\) is \(C^1\) in a neighborhood of \(l = \lim_{p\to\infty} \text{trace}(\Sigma_p)/p = 1\) (for normalized data) and \(C^3\) in a neighborhood of \(0\), satisfying the regularity conditions of~\cite{El_Karoui_2010}.

\paragraph{Step 1: Equivalent Matrix Structure.}
By Theorem 2.1 of~\cite{El_Karoui_2010}, in probability as \(n, p \to \infty\), the kernel matrix \(\mathbf{M}\) is equivalent in operator norm to the matrix \(\mathbf{K}\):
\[
\mathbf{K} = \underbrace{\alpha \mathbf{1}\mathbf{1}^T}_{\text{Spike}} + \underbrace{\beta \frac{XX^T}{p}}_{\text{Bulk}} + \underbrace{\gamma I_n}_{\text{Shift}},
\]
where \(X = [x_1, \dots, x_n] \in \mathbb{R}^{p \times n}\) is the data matrix. The three spectral coefficients \(\alpha, \beta, \gamma\) are determined by the kernel function \(f(\rho) = K_\infty(\rho)\) and its derivatives evaluated at \(\rho = 0\) and \(\rho = 1\):
\begin{align}
\alpha &= f(0) = K_\infty(0), \label{eq:alpha} \\
\beta &= f'(0) = K'_\infty(0), \label{eq:beta} \\
\gamma &= f(1) - f(0) - f'(0) = K_\infty(1) - K_\infty(0) - K'_\infty(0). \label{eq:gamma}
\end{align}

\paragraph{Step 2: Calculation of Spectral Coefficients.}

\textbf{A. The Spike (\(\alpha\)): The Outlier.}
The spike coefficient is the kernel value at orthogonality (\(\rho = 0\)):
\[
\alpha = K_\infty(0) = \Theta^{(1)}(0)\left(1 - \frac{\arccos(0)}{\pi}\right) + 2\Sigma^{(1)}(0) + 1.
\]
Since \(\arccos(0) = \pi/2\), we obtain:
\[
\alpha = \Theta^{(1)}(0)\left(1 - \frac{\pi}{2\pi}\right) + 2\Sigma^{(1)}(0) + 1 = \Theta^{(1)}(0)\left(1 - \frac{1}{2}\right) + 2\Sigma^{(1)}(0) + 1 = \frac{\Theta^{(1)}(0)}{2} + 2\Sigma^{(1)}(0) + 1.
\]
For ReLU activations with normalized inputs, \(\Theta^{(1)}(0) = \Sigma^{(1)}(0) = 0\), hence \(\alpha = 1\). The rank-one mean term \(\alpha \mathbf{1}\mathbf{1}^T\) induces an outlier eigenvalue:
\[
\lambda_{\mathrm{spike}} = n \cdot \alpha = n \cdot K_\infty(0),
\]
of order \(O(n)\), isolated from the bulk spectrum.

\textbf{B. The Bulk (\(\beta\)): The Linear Scale.}
The bulk coefficient is the derivative of the kernel at \(\rho = 0\), capturing the linear sensitivity:
\[
\beta = K'_\infty(0) = \frac{d}{d\rho}\left[\Theta^{(1)}(\rho)\left(1 - \frac{\arccos(\rho)}{\pi}\right) + 2\Sigma^{(1)}(\rho) + 1\right]_{\rho=0}.
\]
Differentiating term by term:
\begin{align*}
\frac{d}{d\rho}\left[\Theta^{(1)}(\rho)\left(1 - \frac{\arccos(\rho)}{\pi}\right)\right] &= \Theta^{(1)'}(\rho)\left(1 - \frac{\arccos(\rho)}{\pi}\right) + \Theta^{(1)}(\rho) \cdot \frac{1}{\pi\sqrt{1-\rho^2}}, \\
\frac{d}{d\rho}\left[2\Sigma^{(1)}(\rho)\right] &= 2\Sigma^{(1)'}(\rho).
\end{align*}
Evaluating at \(\rho = 0\) (where \(\arccos(0) = \pi/2\) and \(\sqrt{1-0} = 1\)):
\[
\beta = \Theta^{(1)'}(0)\left(1 - \frac{1}{2}\right) + \Theta^{(1)}(0) \cdot \frac{1}{\pi} + 2\Sigma^{(1)'}(0) = \frac{\Theta^{(1)'}(0)}{2} + \frac{\Theta^{(1)}(0)}{\pi} + 2\Sigma^{(1)'}(0).
\]
For ReLU kernels with normalized inputs, \(\Theta^{(1)}(0) = 0\) and \(\Theta^{(1)'}(0) = 1/\pi\), \(\Sigma^{(1)'}(0) = 1/(2\pi)\), hence \(\beta = 1/(2\pi) + 1/\pi = 3/(2\pi)\).

The term \(\beta XX^T/p\) has the same spectral distribution as a Wishart matrix scaled by \(\beta\). When \(n \asymp p\) with aspect ratio \(\gamma_{\mathrm{ratio}} = n/r\), the bulk eigenvalues follow a deformed Marchenko–Pastur distribution scaled by \(\beta\), supported on:
\[
\Big[\;\beta(1-\sqrt{\gamma_{\mathrm{ratio}}})^2,\;\; \beta(1+\sqrt{\gamma_{\mathrm{ratio}}})^2\;\Big].
\]

\textbf{C. The Shift (\(\gamma\)): The Diagonal Residual.}
The shift coefficient accounts for the difference between the diagonal value and the linear approximation:
\[
\gamma = K_\infty(1) - K_\infty(0) - K'_\infty(0) = \alpha + \beta \cdot 1 + \text{non-linear terms} - \alpha - \beta = \text{non-linear residual}.
\]
Computing \(K_\infty(1)\) (on-diagonal, where \(\arccos(1) = 0\)):
\[
K_\infty(1) = \Theta^{(1)}(1) \cdot 1 + 2\Sigma^{(1)}(1) + 1 = \Theta^{(1)}(1) + 2\Sigma^{(1)}(1) + 1.
\]
Therefore:
\[
\gamma = \left[\Theta^{(1)}(1) + 2\Sigma^{(1)}(1) + 1\right] - \alpha - \beta.
\]
For normalized ReLU kernels, \(\Theta^{(1)}(1) = \Sigma^{(1)}(1) = 1/2\), so \(K_\infty(1) = 1 + 1 + 1 = 3\), and \(\gamma = 3 - 1 - 3/(2\pi) = 2 - 3/(2\pi) > 0\).

The diagonal shift \(\gamma I_n\) adds \(\gamma\) to all eigenvalues, acting as natural Tikhonov regularization. This ensures the kernel matrix is well-conditioned (bounded away from zero) and stabilizes inversion.

\paragraph{Step 3: Final Spectral Distribution.}
Combining the three components, the spectrum of \(\mathbf{M}\) converges (in probability) to:
\begin{itemize}
\item \textbf{Spike (rank-one outlier):} A single eigenvalue \(\lambda_{\mathrm{spike}} = n \cdot \alpha = n \cdot K_\infty(0)\) of order \(O(n)\), isolated from the bulk.
\item \textbf{Bulk (deformed Marchenko–Pastur):} The remaining \(n-1\) eigenvalues follow a deformed MP distribution, supported on:
\[
\Big[\;\gamma + \beta(1-\sqrt{\gamma_{\mathrm{ratio}}})^2,\;\; \gamma + \beta(1+\sqrt{\gamma_{\mathrm{ratio}}})^2\;\Big],
\]
where \(\gamma_{\mathrm{ratio}} = \lim_{n\to\infty} n/r \in (0,\infty)\). The bulk width is determined by \(\beta = K'_\infty(0)\), and the bulk location is shifted by \(\gamma = K_\infty(1) - K_\infty(0) - K'_\infty(0)\).
\end{itemize}

\paragraph{Interpretation: Fisher–Kibble Decoupling in the Spectrum.}
This spectral decomposition reveals how the microscopic Fisher–Kibble structure manifests macroscopically:
\begin{itemize}
\item \textbf{Fisher's angular fluctuations} (controlled by \(\rho\)) determine the bulk scale \(\beta\) through \(K'_\infty(0)\), governing learning capacity and conditioning.
\item \textbf{Kibble's radial fluctuations} (norm variances) concentrate into the deterministic shift \(\gamma\) through the diagonal term \(K_\infty(1)\), stabilizing the matrix inversion.
\item The \textbf{bias term} in the kernel (the \(+1\)) propagates directly into \(\alpha\), guaranteeing a dominant spike regardless of data structure.
\end{itemize}
The theorem is proved.
\end{proof}
